\documentclass[10pt]{article}
\usepackage[margin=0.85in]{geometry}
\usepackage{amsmath,amssymb,booktabs,tabularx,graphicx,microtype,enumitem,hyperref,xcolor}
\newcolumntype{Y}{>{\raggedright\arraybackslash}X}
\hypersetup{colorlinks=true,linkcolor=blue!45!black,urlcolor=blue!45!black,citecolor=blue!45!black}

\title{Multi-Frame Positional Geometry for Retrieval-Native Attention}
\author{Jorge Sim\~ao\\EInnovator R\&D}
\date{Inception Draft --- August 2026}

\begin{document}
\maketitle

\begin{abstract}
Progressive Retrieval Attention (PRA) retains external resources outside the live sequential prompt and materializes only query-selected native key--value (K/V) memory. Prior work established that external K/V must be transported with a coherent positional frame and that rotary positional encoding (RoPE) permits deferred rebinding. This paper asks a different question: once external memory can be rebound correctly, must independently retrieved resources inhabit one global positional sequence at all? Conventional concatenation assigns arbitrary inter-resource order and query distance even when those relations have no semantic meaning. We study resource-relative positional frames that preserve each resource's internal token geometry while independently controlling its relation to the live query. Holding retrieval, selected K/V, materialization, attention budget, and model weights fixed, we compare globally preserved, globally packed, query-adjacent, rank-distance, score-distance, nearby-band, and randomized positional policies. Primary outcomes include task quality, source-distance robustness, distractor sensitivity, and prediction consistency under irrelevant resource permutations. The result determines whether multi-frame rebinding is merely a clean systems abstraction or should become a calibrated semantic dimension of the PRA runtime.
\end{abstract}

\section{Introduction}
A conventional transformer prompt imposes one total positional order. If independent resources $D_1,D_2,D_3$ are serialized before a query $q$, serialization asserts both an inter-resource order and different resource--query distances. These relationships may be artifacts of packing rather than properties of the task.

PRA changes the setting. Retrieved resources possess stable logical identity outside the live prompt, and native K/V can be rebound at consumption time. A resource $r$ can preserve its internal source geometry,
\[
p'_{r,i}=b_r+p_{r,i},
\]
while choosing $b_r$ independently for each request. Multiple resources may therefore overlap in effective RoPE coordinates without sharing logical identity or physical storage.

The central decomposition is
\[
\boxed{\text{logical identity}}
\neq
\boxed{\text{physical residency}}
\neq
\boxed{\text{positional frame}}.
\]
Paper 1.5 established the machinery needed for correct rebinding. Paper 1.6 evaluates which positional topology retrieval-native memory should use.

\section{From Correct Rebinding to Positional Topology}
For a live local sequence, the query retains an ordinary sequential trajectory:
\[
p_{q,0}=p_{\mathrm{source}}+T_{\mathrm{local}}.
\]
Attention may nevertheless expose
\[
T_{\mathrm{attn}}=T_{\mathrm{selected}}+T_{\mathrm{local}}+T_q
\]
K/V positions. Selected PRA memory increases attention width without becoming additional sequential history.

At a consumer layer,
\[
K=[K_{D_1}^{\phi_1};K_{D_2}^{\phi_2};\cdots;K_{\mathrm{local}};K_q],
\]
with one normal attention softmax. The functions $\phi_r$ alter positional geometry, not resource identity or normalization.

\section{Positional Policies}
We compare seven policies while holding retrieval and selected memory fixed.

\paragraph{P0: Global source.}
Use a correct globally coherent source frame and serve as the Paper-1.5 correctness reference.

\paragraph{P1: Global packed.}
Pack selected resources sequentially into one virtual sequence before local context.

\paragraph{P2: Resource adjacent.}
Each resource independently terminates next to the query. Resource-relative coordinate ranges may overlap.

\paragraph{P3: Rank distance.}
Map retrieval rank to resource--query distance.

\paragraph{P4: Score distance.}
Map a predeclared normalized retrieval score to positional distance.

\paragraph{P5: Nearby non-overlapping bands.}
Place selected resources in bounded nearby bands while preserving one monotonic order.

\paragraph{P6: Random distance.}
Randomize offsets inside a matched distance envelope as a negative control.

For every policy, resource-internal token distances are preserved.

\section{Hypotheses}
\begin{description}[leftmargin=1.7cm,style=nextline]
\item[H1] Incorrect external-memory geometry materially degrades prediction.
\item[H2] Resource-relative frames reduce prediction variance under irrelevant resource permutations.
\item[H3] Query-relative rebinding reduces degradation with original source distance.
\item[H4] A suitable multi-frame policy preserves or improves task quality.
\item[H5] Fully overlapping frames may incur an out-of-distribution penalty because pretrained models learned a single monotonic positional topology.
\end{description}

\section{Permutation Consistency}
For a fixed evidence set $D=\{D_1,\ldots,D_k\}$ we generate irrelevant serialization permutations $\pi_i(D)$. We report accuracy, NLL, target probability, answer-flip rate, and probability variance. A primary consistency measure is
\[
C_{\mathrm{perm}}=
\mathbb{E}_{i,j}
\left[
D_{\mathrm{JS}}
\left(
P(y\mid\pi_i(D),q),
P(y\mid\pi_j(D),q)
\right)
\right].
\]
A compelling positive result is equal or higher quality with lower permutation sensitivity.

\section{Distance, Distractors, and Multiple Resources}
The source-distance experiment moves identical evidence progressively farther from the query under global policies while leaving resource-relative policies unchanged. The distractor experiment sweeps irrelevant selected resources over $\{0,2,4,8,16,32\}$. A multi-resource experiment uses $\{1,2,4,8\}$ relevant resources and includes redundant, complementary multi-hop, and conflicting evidence.

\section{Experimental Methodology}
Primary experiments use the simplest validated native PRA path and \texttt{REFERENCE\_CORRECTNESS} layer/materialization settings. We begin with Qwen3-0.6B and Llama 3.2-1B. Workloads include controlled synthetic facts, HotpotQA/2Wiki/MuSiQue, QASPER, and typed Paper-7 records.

All positional-policy comparisons are paired: selected resources, K/V, token budget, retrieval scores, local prompt, and generation settings are identical. Only positional mapping changes. We report bootstrap confidence intervals and per-model/per-workload effects rather than hiding sign reversals in one aggregate.

\section{Layer-Profile Interaction}
After the reference experiment, a limited factorial subset compares positional policies with calibrated consumer/detail-layer profiles. The question is whether improved positional geometry can preserve quality with fewer native consumer/detail layers, creating a direct quality--memory interaction relevant to Paper 4.5.

\section{Engine Reproduction}
The main experiment should remain engine-independent and run through the simplest validated native path. A small subset is reproduced through one engine-native integration, initially MLX if practical, to test whether observed positional effects survive outside the Hugging Face harness.

\section{Metrics}
Primary outcomes are task accuracy or EM, NLL, target probability, permutation JS/KL divergence, answer-flip rate, calibration where meaningful, distractor sensitivity, and the slope of quality against original source distance. Secondary diagnostics include evidence attention mass, attention entropy, resource-level attention distribution, and logit margin. Systems measurements include rebinding latency and positional-metadata/storage overhead.

\section{Falsification}
The multi-frame quality hypothesis is not supported if a correct global frame matches or beats resource-relative policies on quality and robustness. Negative outcomes include overlapping frames harming accuracy, distractors becoming over-salient, consistency improving only by sacrificing quality, gains being confined to synthetic data, or positional policies failing to transfer across model families.

A negative result remains informative: Paper-1.5 rebinding would then be primarily a correctness and systems abstraction rather than a semantic-quality mechanism.

\section{Implications for the PRA Runtime}
Paper 4.5 should not change defaults before this evidence exists. If no robust advantage emerges, the runtime retains the reference correctness geometry and exposes multi-frame policies only experimentally. If a robust model/workload-specific effect appears, semantic PRA profiles gain an explicit positional-policy field, for example:
\begin{verbatim}
position:
  policy: resource_adjacent
  parameters: {}
  evidence_tier: BENCHMARK
\end{verbatim}
The profile registry must retain the policy and its evidence provenance.

At the engine boundary, three mappings remain independent:
\[
\text{logical resource identity},
\qquad
\text{physical engine placement},
\qquad
\text{request-specific positional mapping}.
\]
vLLM scheduler length, SGLang Radix length, and MLX prompt/rotating-cache length must not be used to infer PRA resource positions.

\section{Related Work}
The final manuscript will position the study relative to RoPE and relative positional encoding, position interpolation and context scaling, lost-in-the-middle effects, retrieval/external-memory attention, permutation sensitivity in long prompts, and structured or relational positional encodings. The review should distinguish prior positional mechanisms from the specific setting studied here: independently retained native-K/V resources with request-specific positional rebinding.

\section{Limitations}
The pretrained models were trained primarily on single monotonic sequences, so overlapping resource-relative frames are out of distribution. Effects may vary by model family, size, fine-tune, workload, and consumer-layer profile. Retrieval quality is intentionally held fixed; this paper does not claim to improve discovery. Engine-specific scheduling and physical-residency performance are secondary to the positional-quality question.

\section{Conclusion}
Correct positional transport is necessary but does not determine the best topology for retrieval-native memory. PRA makes it possible to preserve each resource's internal geometry while choosing its query-relative placement independently. Paper 1.6 tests whether that freedom improves quality, distance robustness, and invariance to arbitrary resource serialization, or whether pretrained models continue to prefer a single global positional sequence. Either result provides a concrete positional-policy recommendation for the PRA runtime and its engine integrations.

\end{document}

